\section*{अध्याय \ref{ch:b3:milky-way-galaxies} --- आकाशगंगा और मंदाकिनियाँ}

\begin{solution}{exo:b3:milky-way-galaxies:1}
(क)~$2\pi R_0 = 2\pi\times\qty{2.523e17}{km} = \qty{1.585e18}{km}$;
\qty{229}{km.s^{-1}} की गति से इसमें \qty{6.92e15}{s} लगते हैं, अर्थात्
\qty{219}{Myr}। (ख)~$4540/219 = 20.7$: लगभग बीस चक्कर।
\end{solution}

\begin{solution}{exo:b3:milky-way-galaxies:2}
(क)~$\qty{1}{au}/\qty{1}{yr} = \num{1.496e8}/\num{3.156e7} =
\qty{4.74}{km.s^{-1}}$। (ख)~$v = 4.74\times6.411\times8.178 =
\qty{248.5}{km.s^{-1}}$, जो समीक्षा के
\qty[separate-uncertainty=true]{248 \pm 3}{km.s^{-1}} मान के भीतर है। (ग)~$248.5 - 11 = \qty{237.5}{km.s^{-1}}$,
लगभग \qty{8}{km.s^{-1}} अधिक है, लाल दानवों के \qty{229}{km.s^{-1}} मान से:
दोनों भिन्न $R_0$ मानों पर आधारित हैं (8.178 और \qty{8.122}{kpc}),
और \omterm{def:b3:milky-way-galaxies:lsr}{स्थानीय विराम मानक} स्वयं लगभग \qty{15}{km.s^{-1}} तक
वास्तविक वृत्तीय कक्षा के सापेक्ष बह सकता है।
\end{solution}

\begin{solution}{exo:b3:milky-way-galaxies:3}
(क)~$n = 10(1 - 0.6) = 4$: E4। (ख)~$b/a = 1 - 7/10 = 0.3$। (ग)~वर्गीकरण
दृश्य अक्ष-अनुपात का उपयोग करता है: अपनी लघु अक्ष की दिशा से देखा गया चपटा दीर्घवृत्तज
वृत्ताकार रूपरेखा रखता है और E0 के रूप में वर्गीकृत होता है।
\end{solution}

\begin{solution}{exo:b3:milky-way-galaxies:4}
$d = 10^{\mu/5 + 1}$~pc: LMC \qty{49.6}{kpc}, एंड्रोमेडा
\qty{783}{kpc}, NGC~4258 \qty{7.58}{Mpc}। मेघ के एक तारे का
लंबन $1/\qty{49.6}{kpc} = \qty{0.020}{mas}$ है, जो सबसे अच्छी सूची की
चमकीले तारों के लिए अनिश्चितता के आकार के बराबर है: एक अकेला लंबन
बिल्कुल भी दूरी नहीं देता, और मेघ की स्थिति अन्य साधनों से तय होती है
(ग्रहणी युग्म तारे), जिनसे सीफ़ीड मापनी शुरू होती है।
\end{solution}

\begin{solution}{exo:b3:milky-way-galaxies:5}
(क)~$\omega_0 = A - B = \qty{27.2}{km.s^{-1}.kpc^{-1}}$ और $V_0 =
27.2\times8.178 = \qty{222}{km.s^{-1}}$। (ख)~$V'_0 = -(A + B) =
\qty{-3.4}{km.s^{-1}.kpc^{-1}}$: धीरे-धीरे नीचे गिरता वक्र। (ग)~दोनों
लाल दानवों के मानों से नीचे हैं (229 और \qty{-1.7}{km.s^{-1}}
प्रति किलोपारसेक); उस वक्र के लेखकों ने भी यही पाया, कि $V_0$ का मान
$R_0$ और इन ऊर्ट स्थिरांकों से लगभग दो मानक विचलन कम है।
ऊर्ट के सूत्र एक अक्षसममित चक्रिका में वृत्ताकार कक्षाएँ मानते हैं; शून्येतर $C$ और
$K$ दिखाते हैं कि स्थानीय वेग-क्षेत्र पट्टी द्वारा
और सर्पिल भुजाओं द्वारा विक्षुब्ध है, और स्थानीय $A$ और $B$ उन विक्षोभों को वहन करते हैं।
\end{solution}

\begin{solution}{exo:b3:milky-way-galaxies:6}
(क)~$\rho = V^2/(4\pi GR^2) = 229^2/(4\pi\times\num{4.301e-6}\times
8.178^2) = \num{1.45e7}$ सौर द्रव्यमान प्रति घन किलोपारसेक, या
\num{0.0145} प्रति घन पारसेक। (ख)~$M_\odot c^2 = \num{1.9884e30}\times
(\num{2.998e8})^2 = \qty{1.787e47}{J} = \qty{1.115e57}{GeV}$, और
$\qty{1}{pc^3} = (\qty{3.0857e18}{cm})^3 = \qty{2.938e55}{cm^3}$:
अनुपात \qty{37.96}{GeV.cm^{-3}} है। फिर $0.30/37.96 = \num{0.0079}$ सौर
द्रव्यमान प्रति घन पारसेक। (ग)~समतल-वक्र घनत्व वह सारा द्रव्यमान गिनता है
जो कक्षाओं को थामे रखता है, जिनमें तारे और गैस भी सम्मिलित हैं, मानो वह एक
गोले में फैला हो; अदीप्त द्रव्य उसका केवल एक भाग है, यहाँ लगभग आधा।
\end{solution}

\begin{solution}{exo:b3:milky-way-galaxies:7}
(क)~$W_I(\mathrm{LMC}) = -3.314\times1.4771 + 15.888 = \qty{10.993}{mag}$,
अतः $M_W = 10.993 - 18.477 = \qty{-7.484}{mag}$। (ख)~$\mu = 25.00 + 7.484 =
32.484$ और $d = 10^{32.484/5 + 1}~\mathrm{pc} = \qty{31.4}{Mpc}$।
(ग)~$\delta d/d = (\ln 10/5)\,\delta\mu = 0.4605\times0.077 =
\qty{3.5}{\percent}$ एक तारे के लिए, $3.5/\sqrt{50} = \qty{0.5}{\percent}$
पचास के लिए, यदि उनकी त्रुटियाँ स्वतंत्र हों (मेघ की अपनी दूरी की
त्रुटि इन सभी में समान रूप से जुड़ती है)।
\end{solution}

\begin{solution}{exo:b3:milky-way-galaxies:8}
$\mu = I - A_I - M_I = 26.00 - 0.05 + 4.05 = 30.00$, अतः $d =
10^{7}~\mathrm{pc} = \qty{10.0}{Mpc}$। \qty{0.04}{mag} की त्रुटि
दूरी में $0.4605\times0.04 = \qty{1.8}{\percent}$ के बराबर है।
\end{solution}

\begin{solution}{exo:b3:milky-way-galaxies:9}
(क)~$d = cz/H_0$ से $5\log_{10}d = 5\log_{10}(cz) - 5\log_{10}H_0$ प्राप्त होता है,
अतः $m_B - M_B - 25 = 5\log_{10}(cz) - 5\log_{10}H_0$;
जहाँ $\log_{10}(cz) = a_B + 0.2\,m_B$, $5\log_{10}H_0 = 5a_B + M_B + 25$, जो
यही परिणाम है। (ख)~$\log_{10}H_0 = -3.8506 + 0.714158 + 5 = 1.8636$, अतः
$H_0 = \qty{73.04}{km.s^{-1}.Mpc^{-1}}$। (ग)~$M_B = 5(\log_{10}67.4 -
5.714158) = \qty{-19.427}{mag}$: अधिनवतारे \qty{0.174}{mag} अधिक चमकीले हैं,
अतः सभी दूरियाँ $10^{0.2\times0.174} = 1.084$ गुना बड़ी हैं,
यह $73.04/67.4$ का अनुपात है। \omterm{def:b3:milky-way-galaxies:tension}{हबल तनाव}
दूरी-पैमाने में आठ प्रतिशत की असहमति है।
\end{solution}

\begin{solution}{exo:b3:milky-way-galaxies:10}
(क)~किनारे पर, वृत्ताकार कक्षा में स्थित एक तारे के लिए $V^2/R \approx GM/R^2$
होता है (गोलीय सन्निकटन)। (ख)~$R = \sqrt{M/(\pi\Sigma)}$, अतः $V^2 =
GM/R = G\sqrt{\pi\Sigma}\,M^{1/2}$ और $M = V^4/(\pi\Sigma G^2) \propto
V^4$; समान द्रव्यमान-दीप्ति अनुपात मानने पर $L \propto V^4$ भी।
(ग)~$(100/229)^4 = \num{0.036}$: तेज़ मंदाकिनी
27 गुना अधिक द्रव्यमान वाली है। (घ)~समान दीप्ति पर फ़्लक्स $1/d^2$ के अनुसार घटता है, अतः कम दीप्त
मंदाकिनी दोगुनी दूरी पर है।
\end{solution}

\begin{solution}{exo:b3:milky-way-galaxies:11}
(क)~$\qty{1}{km.s^{-1}} = \qty{1.0227}{kpc.Gyr^{-1}}$, अतः
\qty{229}{km.s^{-1}} के बराबर \qty{234.2}{kpc.Gyr^{-1}} है। (ख)~$2\pi R/V$:
\qty{4}{kpc} पर \qty{107}{Myr}, \qty{12}{kpc} पर \qty{322}{Myr}।
(ग)~$\Delta\phi = Vt\,(1/4 - 1/12) = 234.2\times0.1667 = \qty{39.0}{rad}$,
6.2 चक्कर: भुजा 4 और
\qty{12}{kpc} के बीच छह कसे चक्करों की कुंडली होगी, जो देखी गई भुजाओं से कहीं अधिक कसी है, और हर अगले
घूर्णन पर और भी कसी होती जाएगी। (घ)~घनत्व-तरंग में प्रतिरूप
एक ही कोणीय वेग से, ठोस पिंड की भाँति घूमता है; तारे और गैस उसमें से होकर गुज़रते हैं, इसलिए
भुजाओं का आकार नहीं लिपटता।
\end{solution}

\begin{solution}{exo:b3:milky-way-galaxies:12}
(क)~$a = (4.2\times10^6\times16^2)^{1/3} = \qty{1024}{au}$;
\qty{8178}{pc} पर यह $1024/8178 = \ang{;;0.125}$ कोण अंतरित करता है
(एक पारसेक पर एक au एक आर्कसेकंड के बराबर होता है)। (ख)~$R_S = 2GM/c^2 = 2\times\num{6.674e-11}
\times\num{8.35e36}/(\num{2.998e8})^2 = \qty{1.24e10}{m} = \qty{0.083}{au}$;
व्यास $2R_S$ का कोण $0.166/8178 = \num{2.0e-5}$ आर्कसेकंड अंतरित करता है, लगभग
\qty{20}{\micro as}। (ग)~$R_S = \qty{4.02e-7}{pc}$, और $M/(\tfrac43\pi
R_S^3) = \num{1.5e25}$ सौर द्रव्यमान प्रति घन पारसेक, जो
सैजिटेरियस~A* की गति और आकार से तय निचली सीमा से चालीस गुना है:
यह सीमा कृष्ण विवर के अपने घनत्व के चालीस गुने के भीतर आती है।
\end{solution}

\begin{solution}{pb:b3:milky-way-galaxies:1}

\textbf{भाग I.}

\textbf{1.} $v_r = A\,d\sin2\ell$: \ang{45} पर $+A d$, \ang{90} पर 0, \ang{135} पर
$-A d$, और \ang{180} पर 0।

\textbf{2.} $A = (15.4 + 15.2)/(2\times\qty{1.00}{kpc}) =
\qty{15.3}{km.s^{-1}.kpc^{-1}}$.

\textbf{3.} \ang{90} पर, $v_t = d(B - A) = 4.74\times(-5.91)\times1.00 =
\qty{-28.02}{km.s^{-1}}$, अतः $B = 15.3 - 28.02 =
\qty{-12.72}{km.s^{-1}.kpc^{-1}}$.

\textbf{4.} \ang{0} पर, $v_t = d(A + B) = 4.74\times0.55 =
\qty{2.61}{km.s^{-1}}$, अतः $B = 2.61 - 15.3 = \num{-12.69}$: दोनों
\num{0.03} के भीतर मेल खाते हैं।

\textbf{5.} $\omega_0 = A - B = \qty{28.0}{km.s^{-1}.kpc^{-1}}$ और
$V_0 = 28.0\times8.18 = \qty{229}{km.s^{-1}}$.

\textbf{6.} $V'_0 = -(A + B) = \qty{-2.6}{km.s^{-1}.kpc^{-1}}$: मंद गति से घटता हुआ।

\textbf{7.} $P = 2\pi/\omega_0$: $2\pi\times\qty{8.18}{kpc}/\qty{229.2}{km.s^{-1}}
= \qty{219}{Myr}$; $4540/219 = 20.7$ चक्कर।

\textbf{भाग II.}

\textbf{8.} स्पर्श बिंदु के लिए ऐसी दृष्टि-रेखा चाहिए जो केंद्र के
सूर्य की तुलना में अधिक निकट से गुज़रे, अर्थात $\sin\ell > 0$ और $\ell < \ang{90}$ (या
प्रतिबिंब चतुर्थांश, $\ang{270} < \ell < \ang{360}$, जहाँ सबसे बड़ा
वेग सबसे ऋणात्मक हो)। सौर वृत्त के बाहर
केंद्र तक की दूरी दृष्टि-रेखा के पूरे मार्ग में बढ़ती जाती है: कोई बिंदु
अलग नहीं किया जाता, और केवल वेग से यह नहीं बताया जा सकता कि गैस कहाँ है।

\textbf{9.} $R = 8.18\sin\ell$: \qtylist{5.26;6.27;7.08;7.69}{kpc}।

\textbf{10.} $V = v_\mathrm{max} + V_0\sin\ell$: $82.9 + 147.2 = 230.1$,
$55.6 + 175.4 = 231.0$, $32.4 + 198.3 = 230.7$, $14.9 + 215.2 =
\qty{230.1}{km.s^{-1}}$.

\textbf{11.} $\Delta\nu = \nu v/c = \num{1.420406e9}\times82.9/299792 =
\qty{393}{kHz}$, कम आवृत्तियों की ओर: गैस दूर जा रही है।

\textbf{12.} लाल दानवों के वेग से लगभग \qty{3}{km.s^{-1}} के भीतर: 5 और
\qty{8}{kpc} के बीच वक्र समतल है।

\textbf{13.} बिंदु-द्रव्यमान के लिए $V = V_0\sqrt{R_0/R} = 229\times\sqrt{8.18/5.26} =
\qty{286}{km.s^{-1}}$ पहले स्पर्श बिंदु पर प्राप्त होता है, जो
पाए गए \qty{230}{km.s^{-1}} से कहीं अधिक तेज़ है: द्रव्यमान केंद्र पर संकेंद्रित नहीं है,
इसका अधिकांश भाग 5 और \qty{8}{kpc} के बीच तथा उससे आगे स्थित है।

\textbf{भाग III.}

\textbf{14.} $\num{1.067e10} + 2\times\num{3.944e10} = \num{8.96e10}$ सौर
द्रव्यमान, समीक्षा के \num[separate-uncertainty=true]{5 \pm 1 e10} से लगभग दोगुना; मॉडल के तारे
संभवतः बहुत भारी हैं, जिससे नीचे पाया गया अदीप्त द्रव्य
कम आँका गया होगा।

\textbf{15.} $V_b^2 = GMR^2/(R^2 + b^2)^{3/2} = \num{4.301e-6}\times\num{1.067e10}
\times400/(400.09)^{3/2}$, $V_b = \qty{47.9}{km.s^{-1}}$.

\textbf{16.} पतली: $(a + b)^2 = 30.80$, $V^2 = \num{4.301e-6}\times\num{3.944e10}
\times400/(430.8)^{3/2}$, $V = \qty{87.1}{km.s^{-1}}$। मोटी:
$(a + b)^2 = 11.56$, $V = \qty{90.1}{km.s^{-1}}$.

\textbf{17.} $V_\mathrm{vis} = \sqrt{47.9^2 + 87.1^2 + 90.1^2} =
\qty{134.2}{km.s^{-1}}$.

\textbf{18.} \qty{8.18}{kpc} पर: 74.8, 108.4 और \qty{127.8}{km.s^{-1}},
कुल मिलाकर \qty{183.5}{km.s^{-1}}, $V_0 = \qty{229}{km.s^{-1}}$ की तुलना में:
दृश्य द्रव्यमान सूर्य की स्थिति पर ही कम पड़ जाता है।

\textbf{19.} $\sqrt{G\times\num{8.96e10}/20} = \qty{138.8}{km.s^{-1}}$। मॉडल का
\qty{134.2}{km.s^{-1}} निकट है, क्योंकि \qty{20}{kpc} पर,
चक्रिका की चार से छह मापनीय लंबाई बाहर, लगभग संपूर्ण दृश्य द्रव्यमान स्थित है
कक्षा के भीतर: दृश्य वक्र केप्लरीय हो गया है।

\textbf{भाग IV.}

\textbf{20.} $M = 20\times200^2/\num{4.301e-6} = \num{1.860e11}$ सौर द्रव्यमान।

\textbf{21.} $M_\mathrm{vis} = 20\times134.2^2/\num{4.301e-6} = \num{8.37e10}$
सौर द्रव्यमान।

\textbf{22.} $M_\text{अदीप्त} = \num{1.860e11} - \num{8.37e10} = \num{1.02e11}$;
क्योंकि दोनों द्रव्यमान $RV^2/G$ हैं, जो समान $R$ पर लिए गए हैं, अतः अंश है
$1 - (V_\mathrm{vis}/V)^2 = 1 - (134.2/200)^2 = \num{0.550}$.

\textbf{23.} $x = 20/14.8 = 1.351$, $\ln(2.351) - 1.351/2.351 = 0.8549 - 0.5748
= 0.2801$, $M = 4\pi\times\num{1.06e7}\times14.8^3\times0.2801 =
\num{1.21e11}$ सौर द्रव्यमान; अदीप्त अंश
$1.21/(1.21 + 0.837) = \num{0.59}$ है। प्रभामंडल को पूरे वक्र पर फ़िट किया गया है, जो
\qty{20}{kpc} पर इसे पार कर जाता है: वहाँ मापन
चिकने मॉडल (\qty{209}{km.s^{-1}}) से लगभग \qty{10}{km.s^{-1}} नीचे गिरते हैं।

\textbf{24.} $x = 8.18/14.8 = 0.553$, $\rho = \num{1.06e7}/(0.553\times1.553^2)
= \num{7.95e6}$ सौर द्रव्यमान प्रति घन किलोपारसेक, \num{0.0080} प्रति
घन पारसेक, या \qty{0.30}{GeV.cm^{-3}}, वह मान जो
फ़िट के लेखक उद्धृत करते हैं।

\textbf{25.} $1 - 8.37/18.60 = \num{0.550}$: \qty{20}{kpc} के भीतर के द्रव्यमान का अदीप्त-द्रव्य अंश
है:
\textbf{\qty[text-series-to-math=true]{55}{\percent}}.
\end{solution}
